\documentclass[10pt,a4paper]{amsart}
\usepackage[margin=19mm]{geometry}
\usepackage{amsmath,amssymb,mathtools}
\usepackage[scr=rsfs,cal=euler]{mathalfa}
\usepackage{xfrac}
\usepackage[unicode,hidelinks,bookmarks=false]{hyperref}
\newcommand{\ie}{i.e.\ }
\newcommand{\cf}{cf.\ }
\newcommand{\resp}{respectively\ }
\newcommand{\Subsection}{\subsection}
\providecommand{\nobreakdash}{\nobreak\hbox{-}}
\setlength{\emergencystretch}{2em}
\multlinegap0pt
\usepackage{bm}
\usepackage{bbm}
\usepackage{dsfont}
\usepackage{xspace}
\usepackage{cancel}
\usepackage{mathtools}
\usepackage{amsthm}
\makeatletter
\@ifundefined{theorem}{\newtheorem{theorem}{Theorem}}{}
\@ifundefined{lemma}{\newtheorem{lemma}[theorem]{Lemma}}{}
\makeatother
\newcommand{\HH}{\mathcal{H}}
\newcommand{\tw}{\textsc{tw}}
\newcommand{\unique}{\textrm{unique}}
\title[Supports for Outerplanar and Bounded Treewidth Graphs]{Supports for Outerplanar and Bounded Treewidth Graphs}
\author{Rajiv Raman, Karamjeet Singh}
\date{Source: arXiv:2504.05039v3 (CC BY 4.0)}
\begin{document}
\maketitle

\noindent\emph{Source and licence.} Supports for Outerplanar and Bounded Treewidth Graphs. Version arXiv:2504.05039v3 (CC BY 4.0), distributed under the Creative Commons Attribution 4.0 International licence, as recorded in the arXiv licence metadata for this exact version and shown on the abstract page https://arxiv.org/abs/2504.05039v3. Full rights notice: licenses/arxiv-2504.05039-CC-BY-4.0.txt.\par\smallskip

\noindent\textit{Editorial scope.} Counted unit: the Lemma whose source label is \texttt{lem:mvi} in the article (source0.tex lines 801--810, source label \texttt{lem:mvi}), delivered together with the article's own complete original proof of it (source0.tex lines 812--883, one \texttt{proof} environment of 72 lines). No mathematics is written, repaired or re-derived: statement and proof are the article's own text line for line. Required context retained uncounted: lem:contfree (lines 418--428); lem:npsep (lines 555--565); lem:easydual (lines 769--775). Every definition, display and label of those lines is retained unchanged. Omitted from this package and not redistributed: 1--417, 429--554, 566--768, 776--800, 811--811, 884--1855. Nothing inside a retained excerpt is omitted or altered: every retained body is delivered exactly as the article's lines are written, so no omission receipt of any kind accompanies this package. Citations: the retained excerpts carry no citation command at all, so no bibliography entry of the article is transcribed and no reference is redistributed. Reference adaptation: none was needed and none was made; every reference inside the delivered unit resolves inside it, so no unresolved reference or citation is produced. Printed slips kept literal: no printed word, symbol, hypothesis or number of the counted statement or of its proof was altered. Layout and class: the article's own class is \texttt{article} with packages including inputenc, url, color, epsfig, amsmath, amsthm, amssymb, hyperref, enumerate, graphicx, wrapfig, lipsum, svg, thmtools; the wrapper is the lane's pinned \texttt{amsart} editorial wrapper at 10pt on A4, which restates the theorem-like environments the retained text needs and adds the article's own macro definitions that the retained text needs (\texttt{\textbackslash HH}, \texttt{\textbackslash tw}, \texttt{\textbackslash unique}), with extra wrapper packages (bm, bbm, dsfont, xspace, cancel, mathtools). The delivered numbering is the unit's own. Assets: no retained span contains a figure environment, a graphics-inclusion command, a PDF-page inclusion, a table or a TikZ picture; no image file is copied into this package, the package declares no asset dependency and no image file is redistributed with it. Licence: version arXiv:2504.05039v3 of this article is distributed under the Creative Commons Attribution 4.0 International licence, as recorded in the arXiv licence metadata for this exact version and shown on the abstract page https://arxiv.org/abs/2504.05039v3. A full rights notice is delivered at licenses/arxiv-2504.05039-CC-BY-4.0.txt. Characters: every retained excerpt is pure ASCII, so the delivered engine, XeLaTeX with its default Latin Modern text font, reports no missing character.
\begin{lemma}
\label{lem:contfree}
Let $(G,\mathcal{H,K})$ be a graph system and  $\mathcal{H}^*\subseteq\mathcal{H}$ be the containment-free
subfamily of $\mathcal{H}$ consisting of all the maximal elements of $(\mathcal{H},\prec_N)$.
Then,
\begin{enumerate}
    \item If $(G,\mathcal{H^*,K})$ has an outerplanar support, then $(G,\mathcal{H,K})$ also has an outerplanar support.
    \item If $(G,\mathcal{H^*,K})$ has a support of treewidth $t$, then $(G,\mathcal{H,K})$ also has a support of treewidth $t$.
\end{enumerate}

\end{lemma}

\begin{lemma}
\label{lem:npsep}
Let $(G,\mathcal{H})$ be a non-piercing graph system with
a tree decomposition $(T,\mathcal{B})$ of $G$.
Let $e=\{x,y\}\in E(T)$ be an edge in $T$ where $x$ is a child of $y$.
Then for $H, H'\in\mathcal{H}_{A_{xy}}$, the following holds:
$(i)$ If $A_{xy}\cap H=A_{xy}\cap H'$ and $H|_x\subset H'|_x$, then
$H'|_{-x}\subseteq H|_{-x}$. $(ii)$ If $A_{xy}\cap H=A_{xy}\cap H'$, 
and $H|_x$ and $H'|_x$ properly intersect,
then $H|_{-x}=H'|_{-x}$, and $(iii)$ If $H\cap A_{xy}\subset H'\cap A_{xy}$, and $H|_x$ and $H'|_x$ properly intersect, then $H|_{-x}\subseteq H'|_{-x}$.
\end{lemma}

\begin{lemma}
\label{lem:easydual}
For a connected graph system $(G,\mathcal{H})$ and a tree decomposition
$(T,\mathcal{B})$ of width $t$ that is $k$-sparse with respect to $\mathcal{H}$,
Algorithm {\bf $k$-SDS} computes a dual support
$Q^*$ with $\tw(Q^*)\le k$.
\end{lemma}

\begin{lemma}
\label{lem:mvi}
Let $(G,\mathcal{H})$ be a non-piercing graph system, 
and $(T,\mathcal{B})$ a tree decomposition of $G$ of width $t$.
Then, there is a connected graph system $(G,\mathcal{H}')$ s.t.
$(T,\mathcal{B})$ is a $2^{4(t+1)}$-sparse with respect to $\unique(\mathcal{H}')$. If
$Q'$ is a dual support for $(G,\unique(\mathcal{H}'))$, there
is a dual support $Q^*$ for $(G,\mathcal{H}')$ s.t. $\tw(Q^*)=\tw(Q')$
and $Q^*$ is also a dual support for $(G,\mathcal{H})$.
\end{lemma}

\begin{proof}
Let $(T,\mathcal{B})$ be a tree decomposition of width $t$. We assume wlog
that $T$ is a binary tree rooted at $\rho$.
To obtain $\mathcal{H}'$, we
do a post-order edge traversal of $T$ and at each
edge, we do a \emph{pushing}. 

The pushing operation is as follows:
Let $z$ be the parent of a node $x$ in $T$.
Having done the pushing on the edges in $T_x$, we do the
following at $\{x,z\}$: 
For each $\emptyset\neq S\subseteq A_{xz}$, choose an $H\in\mathcal{M}_S$.
$H$ is called the \emph{pusher} for $S$ at $A_{xz}$. 
For each $H'\in\mathcal{H}^=_S$ s.t. $H'|_x\setminus H|_x\neq\emptyset$,
replace $H'$ by $H'|_x\setminus H|_x$ in $\mathcal{H}$. 
We say that $H'$ has been \emph{pushed}
by $H$ at $A_{xz}$. 

Observe that once a subgraph $H'$ is pushed by a pusher $H$
at $A_{xz}$, $H'|_x\setminus H|_x\subseteq G_x$, and since pushing is
done in a post-order edge traversal of $T$, 
$H'$ is pushed at most once.
Further, the subgraphs in
$\HH$ intersecting $A_{xz}$ were non-piercing before pushing and
since $H'|_x\setminus H|_x\neq\emptyset$, by part $(i)$ and $(ii)$ of Lemma~\ref{lem:npsep},
$H'|_{-x}\subseteq H|_{-x}$. Hence,
$H'|_x\setminus H|_x=H'\setminus H$. This implies
$H'|_x\setminus H|_x$ is a connected subgraph of $G$.
It follows that there
is a connecting edge $\{u,v\}\in E(G)$ s.t.
$u\in H'|_x\setminus H|_x$ and $v\in H|_x$ as
$H'|_x\setminus H|_x\neq\emptyset$ and $H'\cap H\neq\emptyset$.

Let $\mathcal{H}'$ be the subgraphs at the end of the algorithm.
Note that by construction $\mathcal{H}'$ is in bijective correspondence
with $\mathcal{H}$.
For $S\subseteq A_{xz}$,
if a subgraph $H'\in\HH^{=}_S$ was not pushed by a pusher $H$ at $A_{xz}$, then
$H'|_x=H|_x$. 
Therefore, the number of distinct subgraphs of $G_x$ in the collection $\mathcal{H}'$ intersecting $A_{xz}$ is less than
$2^{t+1}$. That is, 
\begin{align}
\label{eqn:uniq1}
|\unique(\mathcal{H}'_{A_{xz}}|_x)|<2^{t+1}
\end{align}

The subgraphs in $\unique(\HH')$ intersecting the bag $B_x$ can be 
associated with 4-tuples based on its intersection with $A_{xz}$, with the
adhesion sets between $x$ and its children, or with the bag $B_x$ itself.
From Eqn. (\ref{eqn:uniq1}), it follows that there are at most $2^{4(t+1)}$ distinct subgraphs in $\unique(\HH')$ intersecting $B_x$.

By the arguments above, $(G,\unique(\mathcal{H}'))$ is a connected
graph system that is at most $2^{4(t+1)}$-sparse, and hence by 
Lemma~\ref{lem:easydual}, it has a dual support $Q'$ of treewidth at most
$2^{4(t+1)}$. By Lemma~\ref{lem:contfree},
we can extend $Q'$ to a support $Q^*$ for $(G,\mathcal{H}')$ without
increasing the treewidth.

We now argue that $Q^*$ is a dual support for
$(G,\mathcal{H})$. Consider a vertex $v\in V(G)$. 
The algorithm above ensures that each subgraph is pushed at most once.
Suppose $H'\in\mathcal{H}_v$ was pushed by $H$ so that its modified copy $H''=H'\setminus H$
does not cover $v$. Then, $H\in\mathcal{H}_v$.
Let $e =\{a,b\}$ be the connecting edge between $H$ and $H''$ such that $a\in H$ and $b\in H''$.
Since $(T,\mathcal{B})$
is a valid tree decomposition, there is a bag $B$ containing both $a$ and $b$.
Let $H_1$ be the unique representative of $H''$ in $\unique(\HH')$.
Since Algorithm {\bf $k$-SDS} puts a complete graph on the subgraphs intersecting
$B$, it implies $H$ and $H_1$ are adjacent in $Q'$.
In $Q^*$, we made $H'$ adjacent to $H_1$.
Hence, $Q^*$ is a dual support for $(G,\HH)$.
\end{proof}

\end{document}
